\CatchFileBetweenTags{\SolarAngleVal}{calculations/flavor_geometry.tex}{SolarAngleVal}
\CatchFileBetweenTags{\SolarAngleEq}{calculations/flavor_geometry.tex}{SolarAngleEq}
\CatchFileBetweenTags{\SolarAngleExp}{calculations/flavor_geometry.tex}{SolarAngleExperimentalValue}
\CatchFileBetweenTags{\SolarAngleSig}{calculations/flavor_geometry.tex}{SolarAngleSigma}
\CatchFileBetweenTags{\SolarAngleAccText}{calculations/flavor_geometry.tex}{SolarAngleAccText}

\CatchFileBetweenTags{\AtmosAngleVal}{calculations/flavor_geometry.tex}{AtmosAngleVal}
\CatchFileBetweenTags{\AtmosAngleEq}{calculations/flavor_geometry.tex}{AtmosAngleEq}
\CatchFileBetweenTags{\AtmosAngleExp}{calculations/flavor_geometry.tex}{AtmosAngleExperimentalValue}
\CatchFileBetweenTags{\AtmosAngleAccText}{calculations/flavor_geometry.tex}{AtmosAngleAccText}
\CatchFileBetweenTags{\AtmosAngleSig}{calculations/flavor_geometry.tex}{AtmosAngleSigma}

\CatchFileBetweenTags{\AlphaInvVal}{calculations/constants.tex}{AlphaInvVal}

\CatchFileBetweenTags{\ReactorAngleVal}{calculations/flavor_geometry.tex}{ReactorAngleVal}
\CatchFileBetweenTags{\ReactorAngleEq}{calculations/flavor_geometry.tex}{ReactorAngleEq}
\CatchFileBetweenTags{\ReactorAngleExp}{calculations/flavor_geometry.tex}{ReactorAngleExperimentalValue}
\CatchFileBetweenTags{\ReactorAngleAccText}{calculations/flavor_geometry.tex}{ReactorAngleAccText}
\CatchFileBetweenTags{\ReactorAngleSig}{calculations/flavor_geometry.tex}{ReactorAngleSigma}

\CatchFileBetweenTags{\NuCPVal}{calculations/flavor_geometry.tex}{NuCPVal}
\CatchFileBetweenTags{\NuCPEq}{calculations/flavor_geometry.tex}{NuCPEq}
\CatchFileBetweenTags{\NuCPExp}{calculations/flavor_geometry.tex}{NuCPExperimentalValue}
\CatchFileBetweenTags{\NuCPSig}{calculations/flavor_geometry.tex}{NuCPSigma}
\CatchFileBetweenTags{\NuCPAccText}{calculations/flavor_geometry.tex}{NuCPAccText}

\CatchFileBetweenTags{\UeOneVal}{calculations/flavor_geometry.tex}{Ue1Val}
\CatchFileBetweenTags{\UeTwoVal}{calculations/flavor_geometry.tex}{Ue2Val}
\CatchFileBetweenTags{\UeThrVal}{calculations/flavor_geometry.tex}{Ue3Val}
\CatchFileBetweenTags{\UmuOneVal}{calculations/flavor_geometry.tex}{Umu1Val}
\CatchFileBetweenTags{\UmuTwoVal}{calculations/flavor_geometry.tex}{Umu2Val}
\CatchFileBetweenTags{\UmuThrVal}{calculations/flavor_geometry.tex}{Umu3Val}
\CatchFileBetweenTags{\UtauOneVal}{calculations/flavor_geometry.tex}{Utau1Val}
\CatchFileBetweenTags{\UtauTwoVal}{calculations/flavor_geometry.tex}{Utau2Val}
\CatchFileBetweenTags{\UtauThrVal}{calculations/flavor_geometry.tex}{Utau3Val}

\section{The PMNS Matrix: Neutrinos as Lattice Phonons}

The dramatic disparity between the near-diagonal CKM matrix and the near-maximal PMNS matrix is the primary evidence for the structural duality of the substrate. While quarks are \textbf{Topological Solitons} constrained by the friction of geometric interfaces, neutrinos are \textbf{Lattice Phonons}, delocalized vibrations of the bulk. Because phonons possess \textbf{Symmetry Transparency}, their mixing is a harmonic distribution across the global symmetries of the 4D projection.

\subsection{The Primary Bulk Resonance (Solar Angle, \texorpdfstring{$\theta_{12}$}{theta 12})}
The Solar angle defines how a delocalized vibration distributes its energy across the internal interaction sectors. Quarks, as discrete knots, ``see'' the integer hardware channels ($\sigma/\nu = 5/16$). Neutrinos, as continuous waves, ``see'' the continuous geometry of the $H_4$ (600-cell) projection.

In $H_4$ geometry, the 5-fold symmetry manifests through the Golden Ratio ($\phi$). To minimize Entropic Action, the phonon resonance must align with the \textbf{Holographic Aperture} of the 4D manifold. This aperture is defined by the inverse of the Golden bisection ($2\phi$), representing the unique ratio where the phase-front remains coherent across the projected 5-fold axes:
\begin{equation}
    \sin^2 \theta_{12} = \SolarAngleEq \approx \mathbf{\SolarAngleVal}
\end{equation}
This prediction \SolarAngleAccText This identifies $\theta_{12}$ as the \textbf{fundamental resonant mode} of a wave propagating through an icosahedral (5-fold) substrate.

\subsection{The Maximal Frame Bisection (Atmospheric Angle, \texorpdfstring{$\theta_{23}$}{theta 23})}

The Atmospheric angle represents the mixing between the second and third generational tiers. In the $E_8$ substrate, these tiers are anchored by the topological boundary ($\chi=2$). 

By the Principle of Least Action, a system with a 2-point boundary (a simple bisection) reaches a stable steady-state when it maximizes its \textbf{Shannon Entropy}. For a 2-channel system, this bare geometric prior occurs at an exact $50/50$ probability split ($1/\chi = 0.5$). 

However, this $0.5$ represents the bare symmetry at the unification scale. To be observed as a physical parameter on the spacetime manifold, the neutrino phonon must propagate through the lattice. Because $\theta_{23}$ mixes the Areal tier (Generation 2) and the Volumetric tier (Generation 3), the signal expands from the chiral boundary into the internal interaction bulk. 

As derived for the heavy quark sector, this expansion inherently subjects the signal to the \textbf{Volumetric Bulk Shift}, characterized by the ratio of the circular interface to the total saturated capacity ($\pi / (\nu + \sigma)$). The physical, dressed atmospheric angle is the bare maximal entropy state amplified by this volumetric projection:

\begin{equation}
    \sin^2 \theta_{23} = \AtmosAngleEq = \mathbf{\AtmosAngleVal}
\end{equation}

The prediction \AtmosAngleAccText

What standard physics interprets as an arbitrary, continuously running parameter is revealed to be a fixed geometric consequence of translating a maximal-entropy boundary state through the volumetric bulk of the $E_8$ lattice.

\subsection{The Interaction Cross-Talk (Reactor Angle, \texorpdfstring{$\theta_{13}$}{theta 13})}

The Reactor angle represents the unavoidable scattering of bulk phonons off the internal ``color'' pillars of the lattice. Although neutrinos are colorless, they propagate through the same $\nu=16$ channel capacity that hosts the strong force. 

The probability of this scattering is the \textbf{Geometric Admittance} of this color-sector cross-talk. It is the ratio of the \textbf{Interaction Remainder} (the $\sigma-\chi=3$ degrees of freedom that constitute the color sector) to the total \textbf{Vacuum Impedance} ($\alpha^{-1}$):
\begin{equation}
    \sin^2 \theta_{13} = \ReactorAngleEq \approx \mathbf{\ReactorAngleVal}
\end{equation}

This derivation \ReactorAngleAccText. It suggests that the ``3'' in the PMNS matrix is the exact same structural ``3'' that defines the number of colors and generations.

\subsection{The Continuous Phase Lag (Neutrino CP Phase, \texorpdfstring{$\delta_{CP}^\nu$}{delta cpnu})}

While the Quark CP phase is a discrete offset driven by the boundary lock ($\arctan(\sigma/\chi)$), the neutrino phase is the continuous counterpart. It represents the \textbf{Geometric Rotation} required for a bulk wave to complete a cycle within a 5-fold periodic system. 

The unique solution that preserves circular phase-coherence in a $\phi$-governed bulk is the Golden Angle—the proportion of the circle ($360^\circ$) divided by $\phi$:
\begin{equation}
    \delta_{CP}^\nu = \NuCPEq \approx \mathbf{\NuCPVal}
\end{equation}
This prediction \NuCPAccText It distinguishes the $E_8$ lattice from standard theories by providing a fixed geometric origin for leptonic CP violation, rather than a range of possibilities.

\subsection{The Full PMNS Matrix Synthesis: State Saturation}

Just as the unitarity of the CKM matrix is structurally enforced by the finite $\nu=16$ channel capacity of the lattice, the PMNS matrix must also satisfy perfect State Saturation. The sum of the resonance probabilities across the global symmetries cannot exceed the active capacity of the local node. Probability cannot be lost; it is perfectly distributed across the available bulk channels.

By inserting ($\theta_{12}, \theta_{23}, \theta_{13}, \delta_{CP}^\nu$) into the standard unitary parameterization, we generate the complete PMNS matrix:

\begin{equation}
|U_{PMNS}| =
\begin{pmatrix}
\mathbf{\UeOneVal} & \mathbf{\UeTwoVal} & \mathbf{\UeThrVal} \\
\mathbf{\UmuOneVal} & \mathbf{\UmuTwoVal} & \mathbf{\UmuThrVal} \\
\mathbf{\UtauOneVal} & \mathbf{\UtauTwoVal} & \mathbf{\UtauThrVal}
\end{pmatrix}
\end{equation}

This matrix provides the visual and mathematical confirmation of the \textbf{Structural Duality} between Quarks and Leptons. 
\begin{itemize}
    \item \textbf{The CKM Matrix} ($|V_{ud}| \approx 0.974$) is heavily identity-dominant. Topological knots are anchored to the boundary by Manifold Drag ($\pi D$); transitioning away from this identity state incurs massive geometric resistance, resulting in heavily suppressed off-diagonal leakage.
    \item \textbf{The PMNS Matrix} ($|U_{e1}| \approx 0.822$, $|U_{\mu3}| \approx 0.749$) is widely distributed. Lattice Phonons lack topological boundary locks and thus experience no Manifold Drag. Their amplitudes represent free, unsuppressed harmonic resonance across the continuous internal symmetries of the $E_8 \to 4D$ projection.
\end{itemize}


\begin{table*}[h]
\centering
\caption{\textbf{Derivation of PMNS Matrix Elements.}}
\label{tab:pmns_derivation}
\renewcommand{\arraystretch}{1.8} 
\setlength{\tabcolsep}{8pt}
\begin{tabular}{@{} l l l c c r @{}} 
\toprule
\textbf{Mechanism} & \textbf{Symbol} & \textbf{Geometric Logic} & \textbf{Predicted} & \textbf{Target (NuFIT)} & \textbf{Dev.} \\
\midrule
Bulk Resonance & $\sin^2 \theta_{12}$ & Golden Ratio Aperture ($\SolarAngleEq$) & \textbf{\SolarAngleVal} & \SolarAngleExp & $+\SolarAngleSig\sigma$ \\
Maximal Entropy & $\sin^2 \theta_{23}$ & Boundary Limit ($\AtmosAngleEq$) & \textbf{\AtmosAngleVal} & \AtmosAngleExp & $-\AtmosAngleSig\sigma$ \\
Interaction Leak & $\sin^2 \theta_{13}$ & Color-Vacuum Cross-talk ($\ReactorAngleEq$) & \textbf{\ReactorAngleVal} & \ReactorAngleExp & $-\ReactorAngleSig\sigma$ \\
Geometric Rotation & $\delta_{CP}^\nu$ & The Golden Angle ($\NuCPEq$) & \textbf{\NuCPVal} & \NuCPExp & $-\NuCPSig\sigma$ \\
\bottomrule
\end{tabular}
\end{table*}